%%=============================================================================
%% 图表、代码和浮动体
%%=============================================================================

%% 复杂表格
\RequirePackage{threeparttable}
\RequirePackage{dcolumn}
\RequirePackage{multirow}
\RequirePackage{booktabs}
\newcolumntype{d}[1]{D{.}{.}{#1}}% or D{.}{,}{#1} or D{.}{\cdot}{#1}

%% 算法环境
\IfFileExists{algorithm.sty}{
\RequirePackage{algorithm}
\RequirePackage{algorithmic}
\floatname{algorithm}{算法}
\renewcommand{\algorithmicrequire}{\textbf{输入:}}
\renewcommand{\algorithmicensure}{\textbf{输出:}}
}{%
\PackageWarning{xmu-thesis-grd}{algorithm package not found. Please install texlive-science}%
}

%% graphics packages
\RequirePackage{graphicx}

%% 并列子图
\RequirePackage{subcaption}
\RequirePackage{wrapfig}

%% 设置图表相关选项
\RequirePackage{float}
\RequirePackage{caption}
\DeclareCaptionLabelSeparator{zhspace}{~~}
\DeclareCaptionFont{captionfont}{\bfseries \songti \zihao{-4}}
\DeclareCaptionFont{subcaptionfont}{\songti \zihao{-4}}
\captionsetup[sub]{
	font = {subcaptionfont,stretch=1}
}
\captionsetup{
	font = {captionfont,stretch=1.3},
	labelsep = zhspace,
}
\captionsetup[table]{
	position = top,
	aboveskip = 9pt,
	belowskip = 9pt,
}
\numberwithin{table}{chapter}
\captionsetup[figure]{
	position = bottom,
	aboveskip = 5pt,
	belowskip = 0pt,
}
\setlength{\textfloatsep}{15pt}

%% 带来源的图注
\newcommand*{\captionwithsource}[2]{%
	\caption[{#1}]{#1}%
	\vspace{6pt}
	\parbox{\textwidth}{\setstretch{1.2}\zihao{-5}\hspace{2em}#2}
}

%% 如果插入的图片没有指定扩展名，那么依次搜索下面的扩展名所对应的文件
\DeclareGraphicsExtensions{.pdf,.eps,.png,.jpg,.jpeg}

%% 颜色宏包
\RequirePackage{xcolor}

%% 设置源代码显示格式
\RequirePackage{listings}
\lstset{tabsize=4, %
	frame = shadowbox, %把代码用带有阴影的框圈起来
	commentstyle = \color{red!50!green!50!blue!50},%浅灰色的注释
	rulesepcolor = \color{red!20!green!20!blue!20},%代码块边框为淡青色
	keywordstyle = \color{blue!90}\bfseries, %代码关键字的颜色为蓝色，粗体
	showstringspaces = false,%不显示代码字符串中间的空格标记
	stringstyle = \ttfamily, % 代码字符串的特殊格式
	keepspaces = true, %
	breakindent = 22pt, %
	numbers = left,%左侧显示行号
	stepnumber = 1,%
	numberstyle = \tiny, %行号字体用小号
	basicstyle = \footnotesize, %
	showspaces = false, %
	flexiblecolumns = true, %
	breaklines = true, %对过长的代码自动换行
	breakautoindent = true,%
	breakindent = 4em, %
	aboveskip = 1em, %代码块边框
	%% added by http://bbs.ctex.org/viewthread.php?tid=53451
	fontadjust, %
	captionpos = t, %
	framextopmargin = 2pt, %
	framexbottommargin = 2pt, %
	abovecaptionskip = -3pt, %
	belowcaptionskip = 3pt, %
	xleftmargin = 4em, %
	xrightmargin = 4em, % 设定 listing 左右的空白
	texcl=true, %
	% 设定中文冲突，断行，列模式，数学环境输入，listing 数字的样式
	extendedchars=false, %
	columns=flexible, %
	mathescape=true %
	numbersep=-1em %
}
\renewcommand{\lstlistingname}{代码} %% 重命名 Listings 标题头
